% !PRINTED 388
\weberpage{388}
\begin{center}
§ 106. \gesperrt{The Genera of the Ideal Classes.}
\end{center}

In \S\,100 we saw how, from a group of numbers in the field
$\Omega$, we can derive divisions of the ideals into classes, and
how the class number is to be determined for these divisions.

Here we wish to consider the division into \gesperrt{genera} from
this point of view as well.

In order to be able to extend our considerations at once to
 the orders, we first exclude from the region of the \gesperrt{natural
 numbers} all those numbers which have a common divisor with
some arbitrarily assumed number $S$.

Thereupon we take a positive rational modulus $m$, and
include in $S$, among others, all the primes that are contained
 in $m$.

Thereby all numbers are excluded that have a common
 divisor with $m$.

From the remaining numbers and their quotients one can
form a \gesperrt{group of rational fractions}, which we shall denote
by $Z$.

\gesperrt{If we replace (following Gauss) a number $a^{-1}$ by the
integer which, multiplied by $a$, yields a number congruent
to unity, then we can carry over the theory of
congruence with respect to the modulus $m$ without further ado to
the fractional numbers $Z$.}

If then $a/b$ and $a_1/b_1$ are fractions in $Z$, then $a/b \equiv a_1/b_1$,
whenever $ab_1 \equiv a_1 b$ in the ordinary sense.

If, then, we unite all mutually congruent numbers
in $Z$ into one class, $Z$ splits into $\varphi(m)$ classes of numbers, each
of which can be represented by an integral rational number.

If $M$ is the group of the numbers $a$ from $Z$ which satisfy the condition
\[
a \equiv 1 \pmod m
\]
then the classes of numbers (modulo $m$) are the subsidiary groups
to $M$, and we have
\weq{1}{(Z, M) = \varphi(m).}
We now consider, as in \S\,98, the groups $O$ of the integral
and fractional numbers $\omega$ of the quadratic field $\Omega$, \gesperrt{or of
an order $[Q]$ of this field, of all of which we}
%CONT>>
